\documentclass[aspectratio=169,11pt]{beamer}
\makeatletter
\input{../../templates/dune-professional/beamerthemeDUNEProfessional.sty}
\makeatother
\DuneUseWhite
\input{deck-style.tex}
\title{VD Detector Protection System}
\subtitle{Status and work towards a \\ DPS conceptual design}
\author{Manuel Arroyave}
\date{I\&I meeting --- \DpsMeetingDate}
\hypersetup{colorlinks=true,urlcolor=DuneInfo,linkcolor=DuneInfo}
\newcommand{\captionline}[1]{\par{\footnotesize\color{DuneTextMuted}#1\par}}
\input{diagrams/site-boundaries.tex}
\input{diagrams/styles.tex}
\input{meeting-updates.tex}
\hyphenpenalty=10000
\exhyphenpenalty=10000
\begin{document}
\begin{frame}[plain]
  \titlepage
\end{frame}

\DpsKnownAndMissing

\begin{frame}{DPS protects equipment; SC operates and records}
\renewcommand{\arraystretch}{1.18}
\begin{tabularx}{\textwidth}{@{}>{\RaggedRight\arraybackslash}p{2.2cm}>{\RaggedRight\arraybackslash}X@{}}
\textbf{DPS} & Detector Protection System: arbitrates protection, applies permits and inhibits, and checks the response.\\
    \textbf{DSCS} & DUNE Detector Slow Controls System (SC): operates and monitors equipment, coordinates requests and keeps the logs.\\
\textbf{DAQ} & Data acquisition: configures readout and records detector data.\\
\textbf{BMS / Cryo / SCADA} & Building management, SCADA Controls and Cryogenics: operate their systems and report their state to SC.\\
\textbf{Consortia} & Define operating limits, protection actions and allowed response times.\\
\textbf{Local\newline interlocks} & Can trip equipment without SC or the central DPS.\\
\end{tabularx}\par\bigskip
\textbf{Protection trips take priority over operating requests.}\par\medskip
SDSTA issues personnel-safety alarms.
\end{frame}

%\DpsFireScope

\begin{frame}{Facility inputs reach DPS through slow controls}
\centering
\dpsdiagram{overview}\par\smallskip
\raggedright\DuneDiagramCaption{Logical interfaces. The dashed SDSTA connection is under study; two-way arrows carry status and requests.}\par\smallskip
\raggedright

OPC UA carries status and requests. DPS uses SNMP for power distribution units (PDUs) and supplies, and EtherCAT for remote I/O.
\end{frame}

\begin{frame}{The FFDG BMS connection needs an API licence}
\begin{columns}[T,onlytextwidth]
\begin{column}{0.51\textwidth}
\centering\siteboundarymap
\end{column}
\begin{column}{0.47\textwidth}
\textbf{Proposed connection}\par
ISD proposes a Metasys API in Batavia to connect the FFDG DUNE FD BMS to SC.\par\bigskip
\textbf{Tested with SBN building data}\par
The API and OPC UA publication have been exercised.\par\bigskip
\textbf{Needed for live data}\par
The Metasys API licence.
\end{column}
\end{columns}\par\bigskip
\raggedright\DuneDiagramCaption{Service boundaries at SURF. The map shows the scope of the systems, not individual controllers.}
\end{frame}

\DpsFfdgUpdate

\begin{frame}{The SDSTA SCADA interface needs a hardware survey}
\begin{columns}[T,onlytextwidth]
\begin{column}{0.51\textwidth}
\centering\siteboundarymap
\end{column}
\begin{column}{0.47\textwidth}
\textbf{Hardware survey underway}\par
Controller, connection and network ports at the Ross Dry.\par\bigskip
\textbf{Protocols under investigation}\par
Modbus TCP or OPC-UA could be options.\par\bigskip
\textbf{Agreement under discussion}\par
Set of signals to use, faults and maintenance.
\end{column}
\end{columns}\par\bigskip
\raggedright\DuneDiagramCaption{SDSTA services extend beyond the detector caverns. The survey will establish the available points.}
\end{frame}

\DpsSdstaUpdate

\begin{frame}{Cryogenics presents one system to DUNE}
The cryogenic controls are in the Central Utility Cavern (CUC).
Their internal divisions appear as one system to DUNE FD DSCS.\par\bigskip
\renewcommand{\arraystretch}{1.3}
\begin{tabularx}{\textwidth}{@{}>{\RaggedRight\arraybackslash}p{3cm}>{\RaggedRight\arraybackslash}X@{}}
\toprule
\textbf{Function} & \textbf{Detector interface}\\\midrule
Process control & Cryogenics operates and protects its own process.\\
Protection inputs & Liquid-argon level, relief-valve state and boiling-risk conditions can affect HV and bias permits.\\
Information path & Cryogenics $\leftrightarrow$ SC $\leftrightarrow$ DPS. Trace each measurement to its actual sensor.\\
Before operation & Agree readiness states, limits and the response to missing or invalid data.\\
\bottomrule
\end{tabularx}\par\medskip
The detector running conditions are still being agreed.
\end{frame}

\begin{frame}{Trace each interlock from its source to the response}
\centering
\begin{tikzpicture}[
 box/.style={draw=DuneTextMuted,fill=none,text=DuneText,align=center,text width=2.2cm,minimum height=1.25cm,inner sep=5pt,font=\small},
 edge/.style={-{Stealth},draw=DuneTextMuted,line width=1.2pt}]
 \node[box] (s) at (-5.5,0) {Condition\\and source};
 \node[box] (i) at (-2.75,0) {Physical input\\and connection};
 \node[box,fill=none] (c) at (0,0) {DPS / local\\decision};
 \node[box,fill=none] (o) at (2.75,0) {Action and\\equipment};
 \node[box] (f) at (5.5,0) {Equipment state\\and feedback};
 \draw[edge] (s)--(i);\draw[edge] (i)--(c);\draw[edge] (c)--(o);\draw[edge] (o)--(f);
 \end{tikzpicture}\par\smallskip
\raggedright\DuneDiagramCaption{Review each interlock from its source and physical input through the decision, action and equipment response.}\par\smallskip
\raggedright

Start with the interlocks proposed by the consortia in their interface documents (ICDs).\par\bigskip
For each one, identify the existing hardware, missing connections and required additions. Check the time from detection to completed shutdown.\par\bigskip
\textbf{Use the rack lists, drawings, parts lists and power calculations.}
\end{frame}

\DpsConnectionCollection

\begin{frame}{A trip removes the permit; recovery needs a new request}
\centering
\begin{tikzpicture}[
 box/.style={draw=DuneTextMuted,fill=none,text=DuneText,align=center,text width=2.65cm,minimum height=1.4cm,inner sep=5pt,font=\small},
 edge/.style={-{Stealth},draw=DuneTextMuted,line width=1.2pt}]
 \node[box] (i) at (-5.1,0) {Conditions\\and SC requests};
 \node[box,fill=none] (l) at (-1.7,0) {DPS PLC\\protection rules};
 \node[box,fill=none] (a) at (1.7,0) {Permits / inhibits\\SNMP or wired I/O};
 \node[box] (f) at (5.1,0) {Equipment state\\and feedback};
 \draw[edge] (i)--(l);\draw[edge] (l)--(a);\draw[edge] (a)--(f);
 \end{tikzpicture}\par\bigskip
\raggedright

\textbf{Allow:} the request is valid, required conditions are met and no trip is active.\par\medskip
\textbf{Trip:} remove the permit, retain the cause and check the equipment response. SC logs the event.\par\medskip
\textbf{Recover:} reset the trip, then request a separate re-enable.\par\bigskip
\raggedright\DuneDiagramCaption{Proposed controller: a commercial programmable logic controller (PLC) based on BDE, located in the CUC. Loss-of-signal behaviour is set for each protected load.}
\end{frame}

\begin{frame}{Build on the Batavia and SD DPS test setups}
\begin{columns}[T,onlytextwidth]
\begin{column}{0.48\textwidth}
\textbf{Batavia --- HEERC G271}\par\bigskip
\textbf{Reported progress}\par\medskip
Initial BDE demonstrator design completed and reviewed with the WIEC team at BNL.\par\bigskip
\textbf{Next steps}\par\medskip
Apply BNL feedback, test fan tachometer readout, then assemble and install the demonstrator.
\end{column}
\begin{column}{0.48\textwidth}
\textbf{SD --- Nevada room}\par\medskip
\texttt{dune-fd-test01}\par\medskip
\textbf{Available}\par\medskip
PDU, smoke sensors, relay and GIZMo.\par\medskip
Ignition, OPC UA bridges and emulators.\par\bigskip
\textbf{Next step}\par\medskip
Expand the setup along the Batavia model.
\end{column}
\end{columns}\par\bigskip
\textbf{Goal:} exercise protection inputs, trips, feedback, reset and recovery.\par\medskip
%\captionline{Batavia: Trevor's reported status, draft slides 15--16. SD: Manuel's update.}
\end{frame}

\DpsDetectorTransition

\begin{frame}{Ten BDE groups use the existing bottom-drift design}
\centering
\begin{tikzpicture}[
 box/.style={draw=DuneTextMuted,fill=none,text=DuneText,align=center,text width=3.1cm,minimum height=1.05cm,inner sep=6pt},
 edge/.style={-{Stealth},draw=DuneTextMuted,line width=1.2pt}]
 \node[box,fill=none] (io) at (-4,0) {Auxiliary rack\\remote I/O};
 \node[box] (lv) at (1.6,1.1) {2 LV racks\\2 PL506 supplies};
 \node[box,fill=none] (wiec) at (1.6,-1.1) {8 warm crates\\8 PTCs + 48 WIBs};
 \draw[edge] (io.north east) -| (-1.7,1.1) -- node[above,font=\small] {inhibits} (lv.west);
 \draw[arrows={Stealth-Stealth},draw=DuneTextMuted,line width=1.2pt] (io.south east) -| (-1.7,-1.1) -- node[below,font=\small] {EtherCAT} (wiec.west);
 \draw[edge] (lv.east) -- (4.35,1.1) -- node[right,font=\small] {48 V} (4.35,-1.1) -- (wiec.east);
 \end{tikzpicture}\par\smallskip
\raggedright\DuneDiagramCaption{One documented BDE group. The supply inhibit and 48 V path are separate from the two-way EtherCAT I/O link.}\par\smallskip
\raggedright

\textbf{Ten groups:} 10 auxiliary racks, 20 low-voltage racks and 20 PL506 supplies; 80 warm interface electronics crates (WIECs), containing 480 interface boards (WIBs).\par\smallskip
Each crate has a power and timing card (PTC). Ten bias racks provide the separate charge-readout-plane and field-cage interfaces.\par\medskip
\textbf{Available:} panel drawings, parts lists, power calculations and Ignition screens.
\end{frame}

\DpsExtendBde

\begin{frame}{A cooling fault removes power from the affected crate}
\centering
\begin{tikzpicture}[
 box/.style={draw=DuneTextMuted,fill=none,text=DuneText,align=center,text width=2.35cm,minimum height=1.05cm,inner sep=5pt,font=\small},
 edge/.style={-{Stealth},draw=DuneTextMuted,line width=1.2pt}]
 \node[box] (fan) at (-5.1,0) {Fan-status input};
 \node[box,fill=none] (plc) at (-1.7,0) {DPS PLC\\cooling rule};
 \node[box,fill=none] (lv) at (1.7,1) {Remote output\\PL506 inhibit};
 \node[box] (wiec) at (5.1,1) {Affected WIEC\\48 V off};
 \node[box,fill=none] (bias) at (1.7,-1) {Remote output\\ISEG inhibit};
 \node[box] (crp) at (5.1,-1) {Associated\\bias off};
 \draw[edge] (fan)--node[above=18pt,font=\small] {EtherCAT}(plc);
 \draw[edge] (plc.north east) |- (lv.west);\draw[edge] (lv)--(wiec);
 \draw[edge,dashed] (plc.south east) |- (bias.west);\draw[edge,dashed] (bias)--(crp);
 \end{tikzpicture}\par\smallskip
\raggedright\DuneDiagramCaption{Solid: documented fan-to-supply inhibit. Dashed: associated bias shutdown, if required by the VD rule.}\par\smallskip
\raggedright

A cooling fault removes power from the affected crate. DPS checks the supply state and voltage; SC records the trip.\par\bigskip
The dependency and shutdown order for associated bias are still being agreed.
\end{frame}

\begin{frame}{SC and DAQ use distinct DAPHNE service paths}
\centering

\begin{adjustbox}{max width=\linewidth}\input{diagrams/daphne-icd-reference.tex}\end{adjustbox}\par
\raggedright\DuneDiagramCaption{Working SC--PDS--DPS ICD, Fig. 1. PDS: 34 LV feeds; cathode devices use PoF (next slide).}\par\smallskip
{\footnotesize\textbf{Proposed SiPM bias inhibit:} DPS $\rightarrow$ SC $\rightarrow$ DAPHNE; SC returns bias status and informs DAQ.\par}
\end{frame}

\begin{frame}{DPS switches PoF power through the PDU}
\centering

\dpsdiagram{pof}\par\smallskip
\raggedright\DuneDiagramCaption{Sensor inputs reach the local PLC. The dashed box-enable path is optional; the PDU supplies AC power.}\par\smallskip
\raggedright

\textbf{80 boxes: 320 voltage + 80 temperature readings = 400 readings.}\par\medskip
DPS inhibits power over fibre (PoF) through the PDU. A separate box-enable signal would provide another control option, if available.
\end{frame}

\begin{frame}{LV supplies power DAPHNE and PDS calibration boxes}
\textbf{Working PDS LV scope: 34 feeds.}\par\medskip
\centering
\dpsdiagram{pds-lv}\par\smallskip
\raggedright\DuneDiagramCaption{Proposed SNMP control/status path. Power arrows group LV feeds; hardware-inhibit wiring remains open.}\par\medskip
\raggedright
\textbf{To complete:} map loads to channels and inhibit groups; agree trip conditions and recovery.\par\medskip
Check supply state, voltage and current after each trip.
\end{frame}

\begin{frame}{Drift HV and bias supplies have separate inhibit paths}
\textbf{One drift-HV supply; ten bias crates with 576 channels.}\par
480 channels serve the charge readout planes (CRPs); 96 serve the field cage.\par\bigskip
\renewcommand{\arraystretch}{1.35}
\begin{tabularx}{\textwidth}{@{}>{\RaggedRight\arraybackslash}p{2.1cm}>{\RaggedRight\arraybackslash}X>{\RaggedRight\arraybackslash}X@{}}
\toprule
\textbf{Equipment} & \textbf{Control and monitoring} & \textbf{Protection}\\\midrule
Heinzinger drift-HV supply & SC / OPC UA through LabVIEW and cRIO; analog control and readback & DPS remote I/O operates the wired interlock\\
ISEG bias supplies & DPS uses SNMP; SC receives status and logs & DPS remote I/O operates wired channel inhibits\\
\bottomrule
\end{tabularx}\par\bigskip
\textbf{Example:} low liquid-argon level removes the drift-HV permit. A drift-HV trip also shuts down the associated field-cage termination supplies.
\end{frame}

\begin{frame}{Trace each protection input to a temperature sensor}
\centering
\begin{tikzpicture}[
 box/.style={draw=DuneTextMuted,fill=none,text=DuneText,align=center,text width=2.5cm,minimum height=0.85cm,inner sep=5pt,font=\small},
 edge/.style={-{Stealth},draw=DuneTextMuted,line width=1.2pt}]
 \node[box] (rtd) at (-5,1) {Temperature sensors\\RTDs};
 \node[box] (crio) at (-1.5,1) {cRIO readout\\LabVIEW};
 \node[box] (fbg) at (-5,-1) {Fibre temperature\\sensors};
 \node[box] (int) at (-1.5,-1) {Optical readout\\interrogator};
 \node[box] (sc) at (2.5,0) {SC / OPC UA\\logs and status};
 \node[box,fill=none,text width=1.5cm] (dps) at (5.8,0) {DPS};
 \draw[edge] (rtd)--(crio);\draw[edge] (fbg)--(int);
 \draw[edge] (crio.east)-| (sc.north);\draw[edge] (int.east)-| (sc.south);
 \draw[edge] (sc)--(dps);
 \end{tikzpicture}\par\smallskip
\raggedright\DuneDiagramCaption{RTD and fibre-sensor measurements reach SC through separate readout systems. Logical interfaces only.}\par\smallskip
\raggedright

Calibration and Cryogenic Instrumentation (CalCI) measures conditions inside the detector. Temperature profiles guide cooldown and operation.\par\bigskip
For protection, each measurement needs a location, a limit, an allowed response time and a named piece of equipment to shut down.
\end{frame}

\begin{frame}{Proposed laser rules block conflicting PDS requests}
\centering
\dpsdiagram{lasers}\par\smallskip
\raggedright\DuneDiagramCaption{SC exchanges requests and status with DPS, DAQ and the CIB. The local laser-safety path remains independent.}\par\smallskip
\raggedright

\textbf{Four laser stations.} Proposed rule: block conflicting laser and PDS requests. Restore PDS after confirmed emission-off or shutter closure.\par\medskip
Local laser-safety interlocks remain independent.
\end{frame}

\begin{frame}{The purity-monitor inhibit rule still needs agreement}
\centering
\begin{tikzpicture}[
 box/.style={draw=DuneTextMuted,fill=none,text=DuneText,align=center,text width=2.45cm,minimum height=1.2cm,inner sep=5pt,font=\small},
 edge/.style={-{Stealth},draw=DuneTextMuted,line width=1.2pt},
 duplex/.style={arrows={Stealth[length=1.6mm,width=1.2mm]-Stealth[length=1.6mm,width=1.2mm]},draw=DuneTextMuted,line width=1.2pt}]
 \node[box] (prm) at (-5.1,0) {Purity monitor\\Bias / lamp};
 \node[box] (read) at (-1.7,0) {Digitizer\\LabVIEW control};
 \node[box] (sc) at (1.7,0) {SC / OPC UA\\Activity and logs};
 \node[box] (daq) at (5.1,0) {DAQ\\activity tags};
 \draw[duplex] (prm)--(read);\draw[duplex] (read)--(sc);\draw[edge] (sc)--(daq);
 \end{tikzpicture}\par\smallskip
\raggedright\DuneDiagramCaption{Two-way links carry monitor control and status; SC sends activity tags to DAQ. The inhibit rule remains open.}\par\smallskip
\raggedright

\textbf{Proposed approach:} keep taking data and mark the intervals when the flash lamp is active.\par\bigskip
\textbf{Decision still needed:} the older PDS interface document requires SiPM bias to be off during purity-monitor activity. Agree the rule before implementing the inhibit.
\end{frame}

\begin{frame}{Use rack sensors and PDUs for shared protection}
\centering
\dpsdiagram{racks}\par\smallskip
\raggedright\DuneDiagramCaption{Sensors feed the PDU; the outgoing line supplies AC power. Two-way links carry control and status.}\par\smallskip
\raggedright

Local PDU rules switch off affected outlets. Humidity is monitored; a humidity trip needs a separate requirement.\par\medskip
\textbf{TDE + DAQ/timing: 336 microTCA crates. TDE uses 20 mezzanine racks.} DPS controls their PDUs; crate-status and external-inhibit interfaces remain possible additions.
\end{frame}

\begin{frame}{Measure the full path from fault to confirmed shutdown}
\centering
\begin{tikzpicture}[
    box/.style={draw=DuneTextMuted,fill=none,text=DuneText,align=center,
      text width=2.25cm,minimum height=1.0cm,inner sep=5pt},
    edge/.style={-{Stealth},draw=DuneTextMuted,line width=1.2pt}]
    \node[box] (in) at (-4.8,0) {Fault input or\\conflicting request};
    \node[box,fill=none] (plc) at (-1.6,0) {DPS PLC\\arbitration};
    \node[box,fill=none] (out) at (1.6,0) {Supply inhibit\\or PDU command};
    \node[box] (state) at (4.8,0) {Equipment state\\and feedback};
    \draw[edge] (in) -- (plc);
    \draw[edge] (plc) -- (out);
    \draw[edge] (out) -- (state);
  \end{tikzpicture}\par\smallskip
\raggedright\DuneDiagramCaption{Measure detection, arbitration, the output action and equipment feedback as one protection path.}\par\smallskip
\raggedright

Use the existing BDE test setup to measure detection through completed shutdown.\par\medskip
Test local trips, lost communications, reset and re-enable. Check that SC logs the cause, decision and final equipment state.\par\bigskip
No critical function requiring less than 1 ms has been identified so far. The consortium review and hardware tests will check this assumption.
\end{frame}

\DpsMilestones

\begin{frame}{Prepare the hardware review}
\renewcommand{\arraystretch}{1.45}
\begin{tabularx}{\textwidth}{@{}>{\RaggedRight\arraybackslash}p{3.4cm}>{\RaggedRight\arraybackslash}X@{}}
\toprule
\textbf{Work before review} & \textbf{Result}\\\midrule
Trace each inhibit & Existing hardware, additions, cable routes and quantities\\
Test representative paths & Measured trip times, recovery and SC logging\\
Prepare purchases & Panel space, power, drawings, parts, licences, costs and spares\\
\bottomrule
\end{tabularx}\par\bigskip
Order after the readiness review; install to match the site schedule.\par\medskip
Agree the facility-interface and full VD response-matrix review schedule with the consortia and LBNF.
\end{frame}

\appendix
\DpsBackupsTransition

\begin{frame}{Agree responsibility and control-engineering support}
\renewcommand{\arraystretch}{1.45}
\begin{tabularx}{\textwidth}{@{}>{\RaggedRight\arraybackslash}p{4.1cm}>{\RaggedRight\arraybackslash}X@{}}
\toprule
\textbf{Proposed responsibility} & \textbf{Work}\\\midrule
Manuel + control engineer & VD hardware design, SURF interfaces and procurement\\
Terri and Trevor & BDE panels, PLC and I/O implementation, bench tests\\
Geoff + Manuel & SC commands, DPS arbitration and recovery\\
Vivian + consortia / DPS & Equipment requirements and interface-document updates\\
\bottomrule
\end{tabularx}\par\bigskip
\textbf{Support needed:} control-engineering time and access to the existing test stands.
\end{frame}

\input{backup-slides.tex}
\end{document}
